\documentclass[varwidth=16cm, border=2mm]{standalone}
\usepackage{amsmath}\usepackage{amssymb}\usepackage{ctex}
\usepackage{tasks}\usepackage{xcolor}\usepackage{graphicx}
\settasks{label=\Alph*., label-width=1.5em, item-indent=2em}
\begin{document}
\textbf{[0.6]} (2024·全国·高三对口高考)已知椭圆 $\frac{x^2}{9} + y^2 = 1$，过左焦点 F 作倾斜角为 $\frac{\pi}{6}$ 的直线交椭圆于 A,B 两点，则弦 AB 的长为\underline{\hspace{1.5cm}}.

\vspace{0.2cm}\par\noindent\textcolor{red}{\textbf{【答案】} 2}
\par\noindent\textcolor{blue}{\textbf{【解析】} 题步骤】\n**步骤 1: 确定椭圆的几何参数。**\n根据椭圆方程 $\frac{x^2}{9} + y^2 = 1$，我们可以得到：\n半长轴 $a^2 = 9 \Rightarrow a = 3$。\n半短轴 $b^2 = 1 \Rightarrow b = 1$。\n焦距 $c^2 = a^2 - b^2 = 9 - 1 = 8 \Rightarrow c = \sqrt{8} = 2\sqrt{2}$。\n因此，椭圆的左焦点 $F$ 的坐标为 $(-c, 0) = (-2\sqrt{2}, 0)$。\n\n**步骤 2: 确定直线 AB 的方程。**\n直线 AB 过点 $F(-2\sqrt{2}, 0)$，且倾斜角为 $\theta = \frac{\pi}{6}$。\n直线的斜率 $k = \tan(\frac{\pi}{6}) = \frac{\sqrt{3}}{3}$。\n根据点斜式方程 $y - y_0 = k(x - x_0)$，直线 AB 的方程为：\n$y - 0 = \frac{\sqrt{3}}{3}(x - (-2\sqrt{2}))$ \n即 $y = \frac{\sqrt{3}}{3}(x + 2\sqrt{2})$。\n将此方程变形为 $x = \sqrt{3}y - 2\sqrt{2}$，以便后续代入椭圆方程。\n\n**步骤 3: 计算弦 AB 的长度。**\n方法一：利用韦达定理和弦长公式。\n将 $x = \sqrt{3}y - 2\sqrt{2}$ 代入椭圆方程 $\frac{x^2}{9} + y^2 = 1$：\n$\frac{(\sqrt{3}y - 2\sqrt{2})^2}{9} + y^2 = 1$ \n两边同乘以 9 得到：\n$(\sqrt{3}y - 2\sqrt{2})^2 + 9y^2 = 9$ \n展开得：\n$3y^2 - 2 \cdot (\sqrt{3}y) \cdot (2\sqrt{2}) + (2\sqrt{2})^2 + 9y^2 = 9$ \n$3y^2 - 4\sqrt{6}y + 8 + 9y^2 = 9$ \n合并同类项，整理成关于 $y$ 的一元二次方程：\n$12y^2 - 4\sqrt{6}y - 1 = 0$ \n设直线 AB 与椭圆的两个交点分别为 $A(x_1, y_1)$ 和 $B(x_2, y_2)$，那么 $y_1$ 和 $y_2$ 是上述方程的两个根。\n根据韦达定理：\n$y_1 + y_2 = -\frac{-4\sqrt{6}}{12} = \frac{4\sqrt{6}}{12} = \frac{\sqrt{6}}{3}$ \n$y_1 y_2 = \frac{-1}{12}$ \n弦长公式为 $|AB| = \sqrt{(x_1 - x_2)^2 + (y_1 - y_2)^2}$。\n由于 $x = \sqrt{3}y - 2\sqrt{2}$，所以 $x_1 - x_2 = (\sqrt{3}y_1 - 2\sqrt{2}) - (\sqrt{3}y_2 - 2\sqrt{2}) = \sqrt{3}(y_1 - y_2)$。\n将此代入弦长公式：\n$|AB| = \sqrt{(\sqrt{3}(y_1 - y_2))^2 + (y_1 - y_2)^2}$ \n$|AB| = \sqrt{3(y_1 - y_2)^2 + (y_1 - y_2)^2}$ \n$|AB| = \sqrt{4(y_1 - y_2)^2} = 2|y_1 - y_2|$ \n接下来计算 $(y_1 - y_2)^2$：\n$(y_1 - y_2)^2 = (y_1 + y_2)^2 - 4y_1 y_2$ \n$(y_1 - y_2)^2 = (\frac{\sqrt{6}}{3})^2 - 4(-\frac{1}{12})$ \n$(y_1 - y_2)^2 = \frac{6}{9} + \frac{4}{12}$ \n$(y_1 - y_2)^2 = \frac{2}{3} + \frac{1}{3} = 1$ \n所以 $|y_1 - y_2| = \sqrt{1} = 1$。\n因此，弦 AB 的长度为 $|AB| = 2 \times 1 = 2$。\n\n方法二：利用椭圆焦弦长公式。\n椭圆的焦弦长公式为 $L = \frac{2ab^2}{a^2 - c^2\cos^2\theta}$，其中 $\theta$ 是焦弦与 $x$ 轴正方向的夹角。\n已知 $a=3$, $b=1$, $c=2\sqrt{2}$，倾斜角 $\theta = \frac{\pi}{6}$。\n所以 $\cos\theta = \cos(\frac{\pi}{6}) = \frac{\sqrt{3}}{2}$。\n将这些值代入公式：\n$|AB| = \frac{2 \cdot 3 \cdot 1^2}{3^2 - (2\sqrt{2})^2 \cdot (\frac{\sqrt{3}}{2})^2}$ \n$|AB| = \frac{6}{9 - 8 \cdot \frac{3}{4}}$ \n$|AB| = \frac{6}{9 - 6}$ \n$|AB| = \frac{6}{3} = 2$。\n两种方法计算结果一致。\n\n【最终答案】\n2}


\end{document}
